\section{评测与基准}

\subsection{评价指标体系}

\begin{itemize}
    \item \textbf{Sample efficiency}：每次 update 平均消耗多少 transitions；
    \item \textbf{Reward saturation}：reward 曲线 plateau 是否稳定；
    \item \textbf{Entropy decay}：entropy 快速下降表示 exploration collapse；
    \item \textbf{Policy divergence}：KL 及梯度方向是否与 critic signal 保持一致。
\end{itemize}

\begin{importantbox}{审慎设定评测窗口}
测试每个 checkpoint 时要使用固定随机种子和 rollout steps，防止 single-run high reward 误导结论。建议把 best-of-$N$ 分数与均值、标准差一起上报。
\end{importantbox}

\subsection{Benchmark 例子}

\begin{itemize}
    \item \textbf{MuJoCo locomotion}：PPO 作为 baseline，SAC 提供更高样本效率；
    \item \textbf{Atari}：A3C/A2C 展示了批量并行能显著缩短收敛；
    \item \textbf{Robotics sim2real}：严格控制 KL、safety filter 与 entropy；
    \item \textbf{多智能体 StarCraft}：centralized critic + decentralized actor。
\end{itemize}

\begin{knowledgebox}{Benchmark 时常犯的错误}
只在 single random seed 下复现 best score，会掩盖 policy collapse 和 training instability。建议至少跑 3-5 条种子，并把 KL 轨迹与 reward curve 一起画。
\end{knowledgebox}

\subsection{本章小结}

评测不只是看 final reward，还要监控 entropy、KL、critic loss 与 sampling ratio，才能判断 Actor-Critic 是否真正稳定。

